(1) 极坐标动能为 \(m(\dot r^2+r^2\dot\phi^2)/2\)，守恒量 \(L=mr^2\dot\phi\)。消去角速度，\(U_{\rm eff}(r)=L^2/(2mr^2)+kr^2/2+\lambda r^4/4\)。由于 \(U_{\rm eff}''=3L^2/(mr^4)+k+3\lambda r^2>0\) 且 \(r\to0,\infty\) 时势趋无穷，存在且只有一个严格稳定的圆轨道半径。
(2) 圆轨道条件 \(U_{\rm eff}'(r_c)=0\) 给出 \(kr_c+\lambda r_c^3=L^2/(mr_c^3)\)，即 \(L^2=m(kr_c^4+\lambda r_c^6)\)；左端的半径函数严格递增，根唯一。\(\Omega=L/(mr_c^2)\) 且 \(m\Omega^2=k+\lambda r_c^2\)。小扰动 \(r=r_c+\eta\) 满足 \(\ddot\eta+\omega_r^2\eta=0\)，其中 \(\omega_r^2=(4k+6\lambda r_c^2)/m\)。
(3) 一个径向周期 \(2\pi/\omega_r\) 内的方位角近似增加 \(\Delta\phi=2\pi\Omega/\omega_r=2\pi\sqrt{(k+\lambda r_c^2)/(4k+6\lambda r_c^2)}\)。此比值小于 \(1/2\)，所以每次近心点之间的轨道并非闭合椭圆，而是多叶轨线；与完整一周相比的进动为 \(\Delta\phi-2\pi<0\)。所谓进动的数值须明确是相邻近心点间隔，不能误套开普勒势的 \(2\pi\) 公式。
(4) 取 \(E_c(L)=U_{\rm eff}(r_c(L),L)\)，求导时有 \(dE_c/dL=\partial_LU_{\rm eff}+(\partial_rU_{\rm eff})dr_c/dL=L/(mr_c^2)=\Omega\)；第二项因驻点为零。\(\lambda\to0\) 时 \(r_c^4=L^2/(mk)\)，\(\Omega^2=k/m\)，\(\omega_r=2\Omega\)，故径向周期中的角增量为 \(\pi\)，与各向同性谐振子封闭椭圆轨道一致。
